\section{Exteriors, flux forms, and conformal changes}
\label{ch:pre:section}

Throughout this part, the spatial dimension is three, $G=c=1$, and the cosmological constant is zero. Data are smooth up to every specified boundary. A subscript on a norm specifies its metric when needed. For a two-sided surface, the designated unit normal $\nu$ points into the exterior toward the chosen end, and
\[
 H=\Div_\Sigma\nu,\qquad \theta_+=H+\tr_\Sigma K.
\]
Thus outward Euclidean spheres have positive mean curvature. The spacetime convention is $K(Y,Z)=\mathbf g(\nabla_Y n,Z)$ for the future unit normal $n$.

\begin{definition}[Charged exterior]\label{ch:pre:charged-exterior}
A charged exterior consists of a connected orientable smooth three-manifold $N$ with nonempty compact smooth boundary $S$, a Riemannian metric $g$ complete with $S$ included, a symmetric covariant tensor $K$, and vector fields $\cE,\cB$. There is exactly one end. In coordinates $x$ on that end, with $r_x=|x|$, for some $q>1/2$,
\begin{equation}\label{ch:pre:falloff}
 g_{ij}-\delta_{ij}=O_2(r_x^{-q}),\quad
 K_{ij}=O_1(r_x^{-1-q}),\quad
 \cE^i,\cB^i=O_1(r_x^{-2}).
\end{equation}
The notation $O_k(r^{-b})$ bounds coordinate derivatives of order $j\le k$ by $Cr^{-b-j}$. Define the total constraint densities by
\begin{equation}\label{ch:pre:constraints}
 16\pi\mu=R_g+(\tr_gK)^2-|K|_g^2,
 \qquad
 8\pi J_i=\nabla^j\bigl(K_{ij}-(\tr_gK)g_{ij}\bigr).
\end{equation}
The functions $\mu$ and $|J|_g$ are integrable. The matter densities are
\begin{equation}\label{ch:pre:matter}
 \mu_m=\mu-\frac{|\cE|_g^2+|\cB|_g^2}{8\pi},
 \qquad
 J_m=J-\frac{(\cE\times\cB)^\flat}{4\pi},
\end{equation}
where the cross product uses the orientation and metric. They satisfy
\begin{equation}\label{ch:pre:charged-dec}
 \mu_m\ge |J_m|_g,\qquad
 \Div_g\cE=\Div_g\cB=0.
\end{equation}
The boundary satisfies $\theta_+(S)\le0$. We define ADM energy and momentum
by the following finite limits \cite{ArnowittDeserMisner1962}:
\begin{align}
 E&=\frac1{16\pi}\lim_{R\to\infty}
 \int_{|x|=R}(\partial_jg_{ij}-\partial_ig_{jj})n_\delta^i\,dA_\delta,
 \label{ch:pre:adm-energy}\\
 P_i&=\frac1{8\pi}\lim_{R\to\infty}
 \int_{|x|=R}\bigl(K_{ij}-(\tr_gK)g_{ij}\bigr)n_\delta^j\,dA_\delta,
 \label{ch:pre:adm-momentum}\\
 Q_E&=\frac1{4\pi}\lim_{R\to\infty}
 \int_{|x|=R}g(\cE,\nu_R)\,dA_g,
 \qquad
 Q_B=\frac1{4\pi}\lim_{R\to\infty}
 \int_{|x|=R}g(\cB,\nu_R)\,dA_g.
 \label{ch:pre:charges}
\end{align}
In the first two formulas the normal and area element are Euclidean; in the last two, $\nu_R$ is the outward $g$-unit normal. No extension of the physical data across $S$ is required.
\end{definition}

\subsection{Filled enclosing cuts}

A smooth enclosing cut $\Gamma$ is a compact embedded surface separating all of $S$ from the distant end. Its inner side contains the obstacle, and bounded complementary pockets are filled. Components of $\Gamma$ may coincide with components of $S$. The frontier so obtained, including any coincident part, is counted in
\begin{equation}\label{ch:pre:area-infimum}
 a_g(S)=\inf_{\Gamma}\Area_g(\Gamma).
\end{equation}
All cuts are oriented toward the end. They need not be trapped or minimal, and an individual component need not separate $S$ from infinity by itself.

\subsection{The charged bound}

\begin{theorem}\label{ch:intro:main}
Let $(N,g,K,\cE,\cB)$ be a charged exterior in the sense of Definition~\ref{ch:pre:charged-exterior}, with inner boundary $S$, ADM energy and momentum $(E,P)$, and total charges $(Q_E,Q_B)$. Then $E>|P|$. Setting
\[
 m=\sqrt{E^2-|P|^2},\qquad
 Q=\sqrt{Q_E^2+Q_B^2},\qquad
 r=\sqrt{a_g(S)/(4\pi)},
\]
we have
\[
 m\ge Q,\qquad r\le m+\sqrt{m^2-Q^2}.
\]
\end{theorem}

Here $a_g(S)$ is the minimum enclosing area, defined as an infimum over all smooth filled enclosing cuts. Neither the inner boundary nor a competing cut is required to be connected. The boundary condition is only $H+\tr_S K\le0$, with the normal pointing toward infinity. The electromagnetic fields are divergence free, and the remaining matter satisfies the dominant energy condition. The statement places no maximality, symmetry, or zero-momentum restriction on the initial data.

\begin{lemma}\label{ch:pre:area-positive}
For every exterior as above, $0<a_g(S)<\infty$. If the infimum is instead defined using filled finite-perimeter inner sets that retain the obstacle and count its coincident frontier, its value is unchanged.
\end{lemma}

\begin{proof}
After forgetting $K$ and the fields, Definition~\ref{ch:pre:charged-exterior}
gives an exterior in Definition~\ref{n3:intro:exterior-class}, with the
same boundary $S$. Lemma~\ref{bridge:bridge:positive-area} therefore gives
$0<a_g(S)<\infty$.

For the perimeter assertion, apply
Lemma~\ref{n3:reduction:area-compactness} to the constant sequence $g_j=g$.
Local convergence is immediate. The falloff in
Equation~\eqref{ch:pre:falloff} gives uniform end bounds
$c\delta\le g\le C\delta$ and $r_x|\partial g|\le C$, and a fixed
large coordinate sphere is a compact enclosing competitor of finite
area. The lemma consequently identifies the smooth and perimeter
infima. Its competitors are the same bounded filled inner sets used here:
they retain the entire obstacle, bounded complementary pockets are
filled, and the full frontier, including contact with $S$, is counted.
Its outward smoothing preserves enclosure and proves the claimed
equality of infima. The auxiliary filling in that perimeter argument
requires no extension of the physical data or energy condition.
\end{proof}

\subsection{The neutral exterior inequality in this notation}

We first record the form of Theorem~\ref{n3:intro:exterior-inequality} used below. The transfer retains arbitrary compact boundary topology and a potentially disconnected boundary.

\begin{proposition}[Neutral exterior inequality]
\label{ch:pre:neutral-exterior}
Let $(N,g,K)$ be a smooth connected orientable one-ended exterior, complete with its nonempty compact smooth boundary $S$ included. Suppose the metric and tensor satisfy the first two bounds in \eqref{ch:pre:falloff}, the densities \eqref{ch:pre:constraints} are integrable as above, the ADM limits \eqref{ch:pre:adm-energy}--\eqref{ch:pre:adm-momentum} are finite, and $\mu\ge|J|_g$ throughout $N$. If $H+\tr_SK\le0$ and $E>|P|$, then
\[
 \sqrt{E^2-|P|^2}\ge\sqrt{a_g(S)/(16\pi)}.
\]
No complete fill-in across $S$ is part of the hypothesis.
\end{proposition}

\begin{proof}
Apply Theorem~\ref{n3:intro:exterior-inequality} with its exterior
$\Omega=N$. Definition~\ref{ch:pre:charged-exterior}, after the field
requirements are omitted, uses the same smoothness, orientation,
one-endedness and completeness convention as
Definition~\ref{n3:intro:exterior-class}. Equations~\eqref{ch:pre:constraints},
\eqref{ch:pre:adm-energy} and~\eqref{ch:pre:adm-momentum} are exactly its
constraint and ADM conventions. The decay, integrability, weak expansion
and future timelike hypotheses in the proposition are precisely those
of that theorem. The class of filled enclosing cuts is also identical:
all boundary components are enclosed, bounded pockets are filled, and
coincident obstacle frontier is counted. Hence its enclosing infimum is
$a_g(S)$, giving the displayed inequality. The theorem is one-sided and
requires no extension across $S$.
\end{proof}

\subsection{Flux forms and a conformal inequality}

Put $X_1=\cE$, $X_2=\cB$, and
\begin{equation}\label{ch:pre:flux-forms}
 \alpha_i=\iota_{X_i}dV_g.
\end{equation}
The divergence equations are $d\alpha_i=0$. Their fluxes across every filled enclosing cut are $4\pi Q_E$ and $4\pi Q_B$. In every change of metric below, retaining a flux form means defining its new vector field by contraction with the new volume form. This preserves closedness and all fluxes without choosing a vector potential.

For $|Y|_g\le1$, define
\begin{align}
 I_g(Y)&=|X_1|_g^2+|X_2|_g^2
       +2\langle X_1\times X_2,Y\rangle_g,
 \label{ch:pre:I}\\
 D_g(Y)&=8\pi\bigl(\mu+J(Y)\bigr)-I_g(Y).
 \label{ch:pre:D}
\end{align}
Then
\begin{equation}\label{ch:pre:margin}
 I_g(Y)\ge0,\qquad
 \inf_{|Y|_g\le1}D_g(Y)
 =8\pi(\mu_m-|J_m|_g)=:M_0.
\end{equation}
The first assertion uses $2|X_1\times X_2|\le|X_1|^2+|X_2|^2$. The second follows by minimizing the linear term $8\pi J_m(Y)$ over the unit ball. Both $I_g$ and $D_g$ are invariant under a constant $SO(2)$ rotation of the pair of forms. The same elementary inequality and the triangle inequality imply
\begin{equation}\label{ch:pre:neutral-margin}
 8\pi(\mu-|J|_g)\ge M_0.
\end{equation}

\begin{lemma}[Conformal charged comparison]\label{ch:pre:conformal}
Let $u\ge1$ be smooth, and set $g'=u^4g$, $K'=u^2K$, retaining the two flux forms. For every $|Y|_g\le1$,
\begin{equation}\label{ch:pre:conformal-estimate}
 u^4D_{g'}(u^{-2}Y)
 \ge D_g(Y)+4u^{-1}\bigl(-\Delta_gu-|K|_g|du|_g\bigr).
\end{equation}
At a fixed oriented surface,
\begin{equation}\label{ch:pre:conformal-boundary}
 \theta_+'=u^{-2}\bigl(\theta_++4\partial_\nu\log u\bigr).
\end{equation}
Here $\Delta_g=\Div_g\nabla$.
\end{lemma}

\begin{proof}
The conformal scalar-curvature formula and the scaling of the tensor terms in \eqref{ch:pre:constraints} give
\[
 16\pi u^4\mu'=16\pi\mu-8u^{-1}\Delta_gu.
\]
For $f=\log u$, the difference of the two connections is
\[
 \nabla'_AB-\nabla_AB
 =2\bigl(df(A)B+df(B)A-g(A,B)\nabla f\bigr).
\]
Taking the divergence of the rescaled trace reversal
$u^2(K-(\tr_gK)g)$ therefore yields
\[
 8\pi u^2J'=8\pi J+4K(\nabla\log u,\cdot).
\]
The new field vectors are $X_i'=u^{-6}X_i$, so their components in orthonormal frames scale by $u^{-4}$ and
$I_{g'}(u^{-2}Y)=u^{-8}I_g(Y)$. Consequently
\[
 u^4D_{g'}(u^{-2}Y)
 =D_g(Y)+(1-u^{-4})I_g(Y)
 +4u^{-1}\bigl(-\Delta_gu+K(\nabla u,Y)\bigr).
\]
Dropping the nonnegative second term and estimating the last tensor contraction proves \eqref{ch:pre:conformal-estimate}. Finally,
$H'=u^{-2}(H+4\partial_\nu\log u)$ and
$\tr_\Sigma^{g'}K'=u^{-2}\tr_\Sigma^gK$, which give \eqref{ch:pre:conformal-boundary}.
\end{proof}
